\subsection{Discussion}
\label{sec:discussion}

\paragraph{Interpretation of main matrix.}
Table~\ref{tab:main-results} should be read as \emph{relative} baseline comparison under a fixed latency model: absolute RTT values encode documented capture, IQA, resampling, and cloud jitter components (\S\ref{sec:setup}), not live hospital network traces.
The key finding is that B2 simultaneously reduces median RTT versus B0 and raises valid-frame rate to unity on ShezhenV3 test frames scored by the deployable Edge-IQA implementation.
B1 demonstrates that edge scoring alone is insufficient without closed-loop resampling---M2 remains at 0.52 because sub-threshold frames still reach the cloud.

\paragraph{Realism of quality injection.}
Because ShezhenV3 lacks human blur labels, we inject Gaussian defocus on half of the test frames, mirroring validation calibration.
This design choice aligns with robotic failure modes (micro-motion between end-effector and oral cavity) but may under-model exposure flicker from auto-gain cameras; Edge-IQA partially addresses exposure via the $E$ term.
Future datasets with expert quality ratings would enable ROC-based $\tau$ selection rather than median splits.

\paragraph{Edge-IQA overlap.}
Validation pairs exhibit score overlap (M5 $\rho{=}0.47$), indicating that Laplacian--exposure fusion alone does not perfectly separate clear versus synthetically blurred real tongue texture.
ROI cropping to COCO tongue boxes and exposure normalization per acquisition session are promising extensions already supported by the ShezhenV3 annotations.
Nevertheless, overlap does not preclude useful gating: routing policies integrate temporal retry dynamics, and JSONL flags provide post-hoc explanations.

\paragraph{Learned IQA placeholder (B3).}
B3 applies a uniform $+0.08$ boost to $Q_{\mathrm{img}}$, reducing retries (M3 from 0.48 to 0.18) while maintaining M2.
We treat this as an ablation hint for train-split fine-tuning rather than a deployable model; the primary Paper~I claim remains interpretable B2 suitable for regulatory audit trails.

\paragraph{Generalization beyond FR3.}
Although auxiliary trials use Franka fake hardware, Edge-IQA and LangGraph are hardware-agnostic: any edge gateway exposing HTTP \texttt{/vision/evaluate} and skill endpoints can replay the same JSONL protocol.
The ShezhenV3 full validation therefore externalizes the algorithmic contribution from the robot platform.

\paragraph{Ethics and data handling.}
All scoring runs offline on a public research corpus; no patient identifiers are written to JSONL logs.
Images remain on local storage (\texttt{tcm\_paths.yaml}); only aggregate statistics enter the manuscript.

\paragraph{Summary of key findings.}
\begin{enumerate}
  \item ShezhenV3-COCO (6{,}719 images) replaces synthetic-only calibration for Paper~I.
  \item Edge-IQA calibrates at $\tau{=}0.498$ with validation medians $0.564/0.431$ (clear/blur).
  \item M5 Spearman $\rho{=}0.47$ on 1{,}144 validation pairs quantifies real-world score overlap.
  \item B2 achieves M2$=1.00$ on 500$\times$3 test frames with real Edge-IQA scores.
  \item B2 reduces M1 p50 from 383.6\,ms (B0) to 341.2\,ms ($\approx$11\% improvement).
  \item B1 matches B0 on M2 ($\approx$0.52), isolating the benefit of closed-loop resampling.
  \item B3's $+0.08$ boost lowers M3 retries but remains an optional learned-IQA placeholder.
  \item M6 Franka track confirms routing trends without altering Table~\ref{tab:main-results}.
\end{enumerate}
